\section{Introduction}
\label{sec:intro}

The democratization of generative artificial intelligence has fundamentally altered the landscape of digital media creation. Today, producing hyper-realistic manipulated imagery---colloquially known as deepfakes---is practically effortless. While these powerful generative technologies hold tremendous potential for legitimate creative, educational, and entertainment applications, they simultaneously pose a massive threat to digital trust. We have seen deepfakes increasingly leveraged for malicious purposes, ranging from the spread of political misinformation and large-scale fraud to defamation and non-consensual content generation \cite{faceforensics}. As the line between authentic and fabricated media continues to blur, the necessity for robust, reliable, and scalable detection mechanisms has never been more urgent.

When we began exploring this problem, we quickly noticed a prevailing trend in the digital forensics community. Most contemporary deepfake detection pipelines approach the problem by directly feeding raw RGB pixels into increasingly massive Convolutional Neural Networks (CNNs) or computationally heavy Vision Transformers (ViTs). The logic here is straightforward: the networks are tasked with learning the complex spatial and semantic inconsistencies that might betray a forgery. However, as generative architectures like diffusion models and advanced GANs continue to refine their visual blending and synthesis capabilities, these semantic "tells"---such as mismatched lighting, weird textures, or impossible geometry---are becoming vanishingly rare. The generators are simply getting too good at mimicking reality. Furthermore, relying on massive models to scan every frame of high-resolution video is computationally prohibitive for real-world deployment.

Faced with these challenges, we asked ourselves if there was a completely different angle from which to attack the problem. Instead of forcing a network to figure out \textit{what} the image portrays in the semantic space, we decided to investigate the digital fingerprint of \textit{how} the image was actually constructed and saved. When a digital image is manipulated---for instance, when a generated face is spliced onto an authentic background---the newly inserted region rarely shares the same JPEG compression history as the original host image. This subtle discrepancy in the compression grid is invisible to the naked eye but acts as a latent smoking gun.

To exploit this, we turned to Error Level Analysis (ELA) \cite{ela2012}. ELA is a forensic technique that intentionally re-compresses an image at a known quality level and analyzes the pixel-wise differences. This process acts as a magnifying glass for compression inconsistencies, transforming a natural-looking image into a stark heatmap where manipulated regions light up due to differing error levels. 

In this paper, we propose and detail a lightweight, highly efficient detection pipeline that bridges the gap between classic digital forensics and modern deep learning. We move away from the computationally heavy RGB semantic analysis and instead feed ELA heatmaps directly into an EfficientNet-B0 \cite{efficientnet} backbone. By doing so, we force the network to become an expert at identifying compression anomalies rather than visual semantics. Throughout our research, we confronted significant domain shifts between natural imagery and the abstract noise generated by ELA. In the following sections, we outline our journey, detailing how we adapted a pre-trained network to understand these compression grids, the challenges we faced with frozen feature extractors, and the fine-tuning strategies that ultimately led to our robust 86.40\% accuracy rate on the FaceForensics++ benchmark.
